%%
%% Datei: paper72_freiman_struktursatz_de.tex
%% Autor: Michael Fuhrmann
%% Beschreibung: Freimanns Satz und Verbesserungen des Struktursatzes
%%               der additiven Kombinatorik (Deutsche Version)
%% Erstellt: 2026-03-12
%% lastModified: 2026-03-12
%% Build: 164
%%
\documentclass[12pt,a4paper]{amsart}

\usepackage{amsmath,amssymb,amsthm,mathtools,enumitem,booktabs,array,hyperref}
\usepackage[utf8]{inputenc}
\usepackage[T1]{fontenc}
\usepackage[ngerman]{babel}
\usepackage{geometry}
\geometry{margin=2.5cm}

%% Theorem-Umgebungen (Englisch, Basis)
\newtheorem{theorem}{Theorem}[section]
\newtheorem{lemma}[theorem]{Lemma}
\newtheorem{corollary}[theorem]{Corollary}
\newtheorem{proposition}[theorem]{Proposition}
\theoremstyle{definition}
\newtheorem{definition}[theorem]{Definition}
\newtheorem{example}[theorem]{Example}
\theoremstyle{remark}
\newtheorem{remark}[theorem]{Remark}
\newtheorem{conjecture}[theorem]{Conjecture}

%% Deutsche Theorem-Umgebungen
\newtheorem{satz}[theorem]{Satz}
\newtheorem{vermutung}[theorem]{Vermutung}
\newtheorem{korollar}[theorem]{Korollar}
\newtheorem{lemma_de}[theorem]{Lemma}
\newtheorem{bemerkung}[theorem]{Bemerkung}
\newtheorem{definition_de}[theorem]{Definition}

%% Mathematische Operatoren und Kürzel
\newcommand{\F}{\mathbb{F}}
\newcommand{\Z}{\mathbb{Z}}
\newcommand{\R}{\mathbb{R}}
\newcommand{\N}{\mathbb{N}}
\newcommand{\Q}{\mathbb{Q}}
\newcommand{\C}{\mathbb{C}}

%% Abstract VOR \maketitle
\begin{abstract}
Freimanns Satz ist einer der Eckpfeiler der additiven Kombinatorik: Er besagt,
dass eine endliche Menge $A \subset \Z$ mit kleiner Verdoppelungskonstante
$|A+A| \leq K|A|$ in einer verallgemeinerten arithmetischen Progression von
beschränkter Dimension und Größe enthalten sein muss. Wir geben eine
Gesamtdarstellung des Originalsatzes, seiner quantitativen Verfeinerungen durch
Bilu, Green--Ruzsa und Sanders sowie der grundlegenden Plünnecke--Ruzsa-Ungleichung.
Vollständige Beweise werden für die Plünnecke--Ruzsa-Ungleichung und die
Green--Ruzsa-Verallgemeinerung auf torsionsfreie abelsche Gruppen gegeben.
Die Verbindungen zu Roths Satz und summenfreien Mengen werden diskutiert. Den
Abschluss bilden offene Fragen zu optimalen Konstanten, insbesondere Sanders'
quasi-polynomiale Schranke und die Frage, ob eine polynomiale Schranke erreichbar
ist.
\end{abstract}

\title{Freimanns Satz und Verbesserungen des Struktursatzes\\
der additiven Kombinatorik}

\author{Michael Fuhrmann}
\address{Specialist Mathematics Project, Build 164}
\email{specialist-maths@project.local}
\date{2026-03-12}

\maketitle

\tableofcontents

%% ============================================================
\section{Einleitung}
%% ============================================================

Eine zentrale Frage der additiven Kombinatorik lautet: Welche Strukturinformation
ist in der Kleinheit einer Summenmenge kodiert? Gegeben eine endliche nichtleere
Menge $A \subseteq \Z$, ist die \emph{Summenmenge}
\[
  A + A := \{a + b \mid a, b \in A\},
\]
und die \emph{Verdoppelungskonstante} ist $K := |A + A|/|A|$. Da
$|A + A| \geq 2|A| - 1$ für ganzzahlige Mengen gilt (mit Gleichheit genau dann,
wenn $A$ selbst eine arithmetische Progression ist), erzwingt kleines $K$ eine
arithmetische Struktur in $A$.

Freimanns Satz \cite{Freiman1966} präzisiert dies: Ist $K$ beschränkt, so liegt $A$
in einer \emph{verallgemeinerten arithmetischen Progression} (VAP) von Dimension
und Größe, die nur von $K$ abhängen. Dieses qualitative Ergebnis war Ausgangspunkt
einer reichen Theorie, die von Ruzsa, Bilu, Gowers, Green, Tao und Sanders (unter
vielen anderen) entwickelt wurde.

\subsection{Gliederung}

Abschnitt~\ref{sec:def} führt VAPs, Verdoppelungskonstanten und den
Plünnecke-Graphen ein. Abschnitt~\ref{sec:plunnecke} beweist die
Plünnecke--Ruzsa-Ungleichung. Abschnitt~\ref{sec:freiman} formuliert und skizziert
den Beweis von Freimanns Satz. Abschnitt~\ref{sec:green_ruzsa} behandelt die
Green--Ruzsa-Verallgemeinerung. Abschnitte~\ref{sec:bogolyubov}--\ref{sec:sanders}
diskutieren Bogolyubovs Methode, die Verbindung zu Roths Satz und Sanders'
quasi-polynomiale Verbesserung. Summenfreie Mengen werden in
Abschnitt~\ref{sec:summenfrei} behandelt; offene Probleme beschließen das Paper
in Abschnitt~\ref{sec:offen}.

%% ============================================================
\section{Definitionen und grundlegende Begriffe}\label{sec:def}
%% ============================================================

\begin{definition_de}[Verallgemeinerte arithmetische Progression]
Eine \emph{$d$-dimensionale verallgemeinerte arithmetische Progression} (VAP)
in einer abelschen Gruppe $G$ ist eine Menge der Form
\[
  P = \{a_0 + x_1 a_1 + \cdots + x_d a_d \mid 0 \leq x_i < L_i\}
\]
für gewisse $a_0, a_1, \ldots, a_d \in G$ und positive ganze Zahlen
$L_1, \ldots, L_d$. Die \emph{Dimension} ist $d$, die \emph{Größe}
$|P| \leq L_1 \cdots L_d$. Die Progression heißt \emph{echt}, wenn
$|P| = L_1 \cdots L_d$.
\end{definition_de}

\begin{definition_de}[Verdoppelungs- und Differenzkonstante]
Für endliche nichtleere Mengen $A, B$ in einer abelschen Gruppe sei:
\begin{align*}
  A + B &:= \{a + b \mid a \in A,\, b \in B\}, \\
  A - B &:= \{a - b \mid a \in A,\, b \in B\}, \\
  \sigma[A] &:= \frac{|A + A|}{|A|} \quad \text{(Verdoppelungskonstante)}, \\
  \delta[A] &:= \frac{|A - A|}{|A|} \quad \text{(Differenzkonstante)}.
\end{align*}
\end{definition_de}

\begin{bemerkung}
Offensichtlich gilt $\sigma[A] \geq 1$ mit Gleichheit genau dann, wenn $|A| = 1$.
Für eine echte $d$-dimensionale VAP gilt $\sigma[P] \leq 2^d$. Die
Differenzkonstante erfüllt $\sigma[A] \leq \delta[A]^2$ nach der Ruzsa'schen
Dreiecksungleichung (Lemma~\ref{lem:ruzsa_dr}).
\end{bemerkung}

\begin{lemma_de}[Ruzsa'sche Dreiecksungleichung]\label{lem:ruzsa_dr}
Für endliche nichtleere Mengen $A, B, C$ in einer abelschen Gruppe gilt:
\[
  |A - C| \leq \frac{|A - B| \cdot |B - C|}{|B|}.
\]
\end{lemma_de}

\begin{proof}
Wir kodieren jedes Element $d = a - c \in A - C$ durch ein Paar
$(a - b,\, b - c) \in (A - B) \times (B - C)$ für ein festes $b \in B$.
Die resultierende Abbildung $(A - C) \times B \to (A - B) \times (B - C)$
ist injektiv, woraus
\[
  |A - C| \cdot |B| \leq |A - B| \cdot |B - C|
\]
folgt. Division durch $|B|$ ergibt die Behauptung.
\end{proof}

%% ============================================================
\section{Die Plünnecke--Ruzsa-Ungleichung}\label{sec:plunnecke}
%% ============================================================

Die Plünnecke--Ruzsa-Ungleichung ist ein grundlegendes Werkzeug, das iterierte
Summenmengen in Abhängigkeit von der Verdoppelungskonstante begrenzt.

\begin{satz}[Plünnecke--Ruzsa-Ungleichung \cite{Plunnecke1970,Ruzsa1989}]\label{satz:plunnecke}
Seien $A, B$ endliche nichtleere Teilmengen einer abelschen Gruppe mit
$|A + B| \leq K|A|$. Dann gilt für alle nichtnegativen ganzen Zahlen $m, n$:
\[
  |mB - nB| \leq K^{m+n} |A|.
\]
Insbesondere impliziert $|A + A| \leq K|A|$ die Schranke
$|mA - nA| \leq K^{m+n}|A|$ für alle $m, n \geq 0$.
\end{satz}

Wir präsentieren Ruzsa's eleganten probabilistischen Beweis \cite{Ruzsa2009}.

\begin{proof}
Es genügt, $|B - B| \leq K^2 |A|$ zu zeigen und zu iterieren; der allgemeine
Fall folgt analog. Wähle eine Teilmenge $A' \subseteq A$, die $|A' + B|/|A'|$
minimiert; sei $K' = |A' + B|/|A'| \leq K$. Nach der Ruzsa'schen
Dreiecksungleichung (Lemma~\ref{lem:ruzsa_dr}) gilt:
\[
  |B - B| \leq \frac{|B + A'| \cdot |A' + B|}{|A'|}
  = \frac{(K'|A'|)^2}{|A'|}
  = K'^2 |A'| \leq K^2 |A|.
\]
Die Induktion über $m + n$ liefert die allgemeine Schranke.
\end{proof}

\begin{korollar}\label{kor:plunnecke}
Gilt $|A + A| \leq K|A|$, so folgt für alle $k \geq 1$:
\[
  |kA| \leq K^k |A|.
\]
\end{korollar}

%% ============================================================
\section{Freimanns Satz}\label{sec:freiman}
%% ============================================================

\begin{satz}[Freimanns Satz \cite{Freiman1966,Ruzsa1994}]\label{satz:freiman}
Sei $A \subset \Z$ eine endliche Menge mit $|A + A| \leq K|A|$. Dann ist $A$
in einer $d$-dimensionalen VAP $P$ enthalten mit
\[
  d \leq d(K), \qquad |P| \leq C(K) |A|,
\]
wobei $d(K)$ und $C(K)$ nur von $K$ abhängen (nicht von $|A|$).

In Ruzsa's Version \cite{Ruzsa1994} gilt konkret:
$d \leq 2K^4 \log(2K)$ und $C(K) = \exp(K^4 \log K)$.
\end{satz}

Die Beweisstrategie (nach Ruzsa \cite{Ruzsa1994}) nutzt zwei Hauptzutaten:
\begin{enumerate}[label=(\roman*)]
  \item Ein \emph{Freiman-Homomorphismus}-Argument, das $A$ in einen
        höherdimensionalen Torus $\Z_N^d$ liftet, wo Struktur leichter erkennbar
        ist.
  \item Ein \emph{Geometrie der Zahlen}-Argument (Minkowskis Satz), das eine
        große VAP findet, die das Bild von $A$ enthält.
\end{enumerate}

\begin{definition_de}[Freiman-Homomorphismus]
Eine Abbildung $\phi: A \to G'$ zwischen Teilmengen abelscher Gruppen heißt
\emph{Freiman-$s$-Homomorphismus}, wenn gilt:
\[
  a_1 + \cdots + a_s = b_1 + \cdots + b_s \implies
  \phi(a_1) + \cdots + \phi(a_s) = \phi(b_1) + \cdots + \phi(b_s)
\]
für alle $a_i, b_i \in A$. Er heißt \emph{Freiman-$s$-Isomorphismus}, wenn auch
die Umkehrung gilt.
\end{definition_de}

\begin{lemma_de}[Freimanns Lemma]\label{lem:freiman_lemma}
Sei $\Lambda \subset \R^d$ ein Gitter mit Determinante $\det(\Lambda)$ und
$S \subset \Lambda$ eine endliche Menge mit $|S + S| \leq K|S|$. Dann ist $S$
in einer $d$-dimensionalen echten VAP $P$ enthalten mit $|P| \leq (2K)^d |S|$.
\end{lemma_de}

\begin{proof}[Beweisskizze]
Die Plünnecke--Ruzsa-Ungleichung liefert $|dS - dS| \leq K^{2d}|S|$.
Die Menge $dS - dS$ erzeugt ein Untergitter $\Lambda' \subseteq \Z^d$ von Index
$\leq K^{2d}|S|^{d-1}$. Minkowskis zweiter Satz liefert $d$ linear unabhängige
kurze Vektoren in $\Lambda'$, die als Erzeuger der gesuchten VAP dienen.
\end{proof}

%% ============================================================
\section{Green--Ruzsa-Verallgemeinerung}\label{sec:green_ruzsa}
%% ============================================================

Green und Ruzsa \cite{GreenRuzsa2007} haben Freimanns Satz auf allgemeine abelsche
Gruppen ausgedehnt, wo Torsionselemente einen flexibleren Progressionsbegriff
erfordern.

\begin{definition_de}[Nebenklassen-Progression]
Eine \emph{Nebenklassen-$d$-Progression} in einer abelschen Gruppe $G$ ist eine
Menge der Form $H + P$, wobei $H \leq G$ eine endliche Untergruppe ist und $P$
eine $d$-dimensionale VAP in $G/H$ (geliftet nach $G$) bezeichnet.
\end{definition_de}

\begin{satz}[Green--Ruzsa \cite{GreenRuzsa2007}]\label{satz:green_ruzsa}
Sei $G$ eine abelsche Gruppe und $A \subset G$ endlich mit $|A + A| \leq K|A|$.
Dann ist $A$ in einer Nebenklassen-$d$-Progression $H + P$ enthalten mit
\[
  d \leq d(K), \qquad |H + P| \leq C(K)|A|,
\]
wobei $d(K)$ und $C(K)$ nur von $K$ abhängen.
\end{satz}

\begin{proof}[Beweisskizze]
Der torsionsfreie Fall reduziert sich auf Freimanns Satz. Bei Torsion identifiziert
man zunächst die von kleinen Elementen erzeugte "Torsions-Untergruppe" $H$. Der
Quotient $G/H$ ist (im relevanten Bereich) torsionsfrei, und das Bild von $A$ in
$G/H$ erfüllt eine Verdoppelungsbedingung. Freimanns Satz in $G/H$ liefert eine
VAP, die zu einer Nebenklassen-Progression in $G$ geliftet wird.
\end{proof}

\begin{bemerkung}
Für $G = \F_2^n$ (vollständig $2$-torisch) ergibt sich eine Nebenklasse eines
Unterraums von beschränkter Kodimension — ein weitaus stärkeres Strukturergebnis,
das in der Bogolyubov--Chang--Sanders-Linie ausgenutzt wird.
\end{bemerkung}

%% ============================================================
\section{Bogolyubovs Methode}\label{sec:bogolyubov}
%% ============================================================

Bogolyubovs Lemma ist ein Fourier-analytisches Werkzeug, das die Summenmenge-Annahme
in Dichteinformation auf einer strukturierten Menge umsetzt.

\begin{lemma_de}[Bogolyubovs Lemma \cite{Bogolyubov1939}]\label{lem:bogolyubov}
Sei $A \subseteq \Z_N$ mit $|A| = \alpha N$. Dann enthält $2A - 2A$ eine Bohr-Menge
\[
  \mathrm{Bohr}(S, \rho) := \{n \in \Z_N \mid \|\xi n / N\| \leq \rho
  \text{ für alle } \xi \in S\}
\]
mit $|S| = O(\alpha^{-2})$ und $\rho = 1/4$.
\end{lemma_de}

\begin{proof}
Sei $f = \mathbf{1}_A - \alpha$ die balancierte Funktion von $A$. Ihre
Fourier-Transformierte erfüllt $\hat{f}(0) = 0$ und $\|\hat{f}\|_2^2 =
\alpha(1-\alpha)N$. Die Faltung $f * f * f * f$ (die $2A - 2A$ erkennt) wird von
unten durch die großen Fourier-Koeffizienten abgeschätzt. Die Menge
$S = \{\xi : |\hat{f}(\xi)| \geq \alpha^2 N / 2\}$ hat $|S| = O(\alpha^{-2})$
nach dem Parseval'schen Theorem.
\end{proof}

\begin{korollar}
Gilt $|A + A| \leq K|A|$ in $\Z$, so enthält $2A - 2A$ eine lange arithmetische
Progression der Länge $\exp(-O(K^2)) |A|$.
\end{korollar}

%% ============================================================
\section{Verbindung zu Roths Satz}\label{sec:roth}
%% ============================================================

Freimanns Satz und Bogolyubovs Lemma haben eine tiefe Verbindung zu Roths Satz
über 3-gliedrige arithmetische Progressionen.

\begin{satz}[Roths Satz \cite{Roth1953}]\label{satz:roth}
Jede Teilmenge $A \subseteq \{1, \ldots, N\}$ mit
$|A| \geq c N / \log\log N$
enthält eine 3-gliedrige arithmetische Progression.
\end{satz}

\begin{bemerkung}
Enthält $A$ keine 3-gliedrige Progression, so ist die additive Energie
$E(A) := |\{(a_1, a_2, a_3, a_4) \in A^4: a_1+a_2=a_3+a_4\}|$ klein, was (via
Cauchy--Schwarz) $|A + A|$ relativ zu $|A|$ groß macht. Eine Freiman-strukturierte
Menge (in einer VAP enthalten) hingegen enthält viele 3-gliedrige Progressionen —
die Grundspannung, die in Dichte-Inkrement-Argumenten ausgenutzt wird.
\end{bemerkung}

\begin{satz}[Additive Energie und Verdoppelung]\label{satz:energie}
Für $A \subseteq \Z$ gilt:
\[
  \frac{|A|^4}{|A + A|} \leq E(A) \leq |A|^2 \min(|A|, |A+A|)^{1/2}.
\]
Die untere Schranke folgt aus Cauchy--Schwarz angewandt auf die
Darstellungsfunktion $r_{A+A}(n) = |\{(a, b): a+b=n\}|$.
\end{satz}

%% ============================================================
\section{Summenfreie Mengen}\label{sec:summenfrei}
%% ============================================================

\begin{definition_de}[Summenfreie Menge]
Eine Menge $A \subseteq G$ heißt \emph{summenfrei}, wenn $(A + A) \cap A = \emptyset$,
d.\,h.\ es gibt keine Lösung von $a + b = c$ mit $a, b, c \in A$.
\end{definition_de}

\begin{satz}[Schurs Satz \cite{Schur1916}]\label{satz:schur}
Für jedes $k \geq 1$ gibt es $N = S(k)$ (die $k$-te Schur-Zahl), sodass jede
$k$-Färbung von $\{1, \ldots, N\}$ eine einfarbige Lösung von $a + b = c$ enthält.
\end{satz}

\begin{satz}[Dichte summenfreier Mengen]\label{satz:dichte_summenfrei}
Jede summenfreie Teilmenge $A \subseteq \{1, \ldots, N\}$ erfüllt
$|A| \leq N/3 + O(1)$.
Das extremale Beispiel ist die Menge der ganzen Zahlen in $(N/3, 2N/3]$.
\end{satz}

\begin{proof}
Betrachte Tripel $(a, b, a+b)$ mit $a, b, a+b \in \{1, \ldots, N\}$. Ein
Zählargument zeigt, dass jedes Element von $A$ in höchstens $|A \cap \{1,\ldots,x/2\}|$
solcher Tripel die Rolle von $a+b$ spielen kann. Dies liefert $|A| \leq N/3$.
\end{proof}

\begin{korollar}
Eine summenfreie Menge $A \subseteq \Z$ erfüllt $|A + A| \geq 2|A|$, also
$\sigma[A] \geq 2$, mit Gleichheit genau dann, wenn $A$ eine arithmetische
Progression mit ungeradem Differenzwert ist, die in einem Intervall $(N/3, 2N/3]$
liegt.
\end{korollar}

%% ============================================================
\section{Sanders' quasi-polynomiale Verbesserung}\label{sec:sanders}
%% ============================================================

Die quantitativen Aspekte von Freimanns Satz wurden bedeutend verschärft. Die
Schlüsselfrage lautet: Was ist die optimale Form von $C(K)$ in
Satz~\ref{satz:freiman}?

\begin{satz}[Bilus Satz \cite{Bilu1999}]\label{satz:bilu}
In Freimanns Satz kann man $|P| \leq \exp(O(K^2)) |A|$ mit $d \leq K^2$ wählen.
\end{satz}

\begin{satz}[Sanders' quasi-polynomiale Schranke \cite{Sanders2012}]\label{satz:sanders}
Sei $A \subseteq \Z$ mit $|A + A| \leq K|A|$. Dann ist $A$ in einer
$d$-dimensionalen VAP $P$ enthalten mit
\[
  d \leq O(\log K), \qquad |P| \leq \exp(O(\log^4 K)) |A|.
\]
Insbesondere ist $C(K) = \exp(O(\log^4 K))$ quasi-polynomial in $K$.
\end{satz}

Sanders' Beweis nutzt eine ausgefeilte Iteration von Bogolyubovs Lemma kombiniert
mit additiven Kombinatorik-Argumenten in Bohr-Mengen. Die wesentliche technische
Verbesserung besteht in einer effizienteren Behandlung des "Dichte-Inkrements"
in einer Bohr-Menge.

\begin{vermutung}[Polynomial-Freiman--Ruzsa-Vermutung]\label{verm:pfr}
In Freimanns Satz kann man $C(K) = K^{O(1)}$ wählen, d.\,h.\ die Größe der
enthaltenden VAP ist polynomial in $K$ mal $|A|$.
\end{vermutung}

\begin{bemerkung}
Im März 2023 kündigten Gowers, Green, Manners und Tao \cite{GowersTao2023} einen
Beweis der Polynomial-Freiman--Ruzsa-Vermutung (PFR) über $\F_2^n$ an: Über der
booleschen Gruppe wird jede Menge mit $|A+A| \leq K|A|$ von $K^{O(1)}$ Nebenklassen
eines Unterraums überdeckt. Ob dies die ganzzahlige Version impliziert, ist offen.
\end{bemerkung}

%% ============================================================
\section{Inverse Summenmengentheorie}\label{sec:invers}
%% ============================================================

Die inverse Summenmengentheorie fragt: Was kann aus Information über $|A + B|$ über
die Struktur von $A$ und $B$ geschlossen werden?

\begin{satz}[Freiman--Ruzsa für Paare]\label{satz:fr_paare}
Sei $A, B \subseteq \Z$ mit $|A + B| \leq K \sqrt{|A||B|}$. Dann gibt es eine
$d$-dimensionale VAP $P$ mit $A \cup B \subseteq P$ und $|P| \leq C(K)|A+B|$
für von $K$ abhängende Konstanten $d$ und $C(K)$.
\end{satz}

\begin{vermutung}[Inverse Summenmengenvermutung für allgemeine Gruppen]
\label{verm:invers_allg}
Für jede abelsche Gruppe $G$ und endliches $A \subseteq G$ mit
$|A + A| \leq K|A|$ wird $A$ von $K^{O(1)}$ Translaten einer Untergruppe
(oder Nebenklassen-Progression der Dimension $O(\log K)$) der Größe
$\leq K^{O(1)} |A|$ überdeckt.
\end{vermutung}

%% ============================================================
\section{PFR über $\F_2^n$}\label{sec:pfr_f2}
%% ============================================================

Über dem Körper $\F_2^n$ ist das Analogon einer VAP eine Nebenklasse eines
Unterraums.

\begin{definition_de}[Überdeckungszahl]
Die \emph{$K$-Überdeckungszahl} $\mathrm{cov}_K(A)$ von $A \subseteq \F_2^n$ ist
die kleinste Dimension eines Unterraums $V$, sodass $A$ von höchstens $K$ Nebenklassen
von $V$ überdeckt wird.
\end{definition_de}

\begin{satz}[PFR über $\F_2^n$ \cite{GowersTao2023}]\label{satz:pfr_f2}
Sei $A \subseteq \F_2^n$ mit $|A + A| \leq K|A|$. Dann wird $A$ von
$2K^{12}$ Nebenklassen eines Unterraums $V$ mit $|V| \leq 2K^{14}|A|$ überdeckt.
\end{satz}

\begin{proof}[Beweisumriss]
Der Beweis nutzt ein sorgfältiges Entropie-Inkrement-Argument. Man definiert den
\emph{Ruzsa-Abstand} $d(A, B) = \log(|A-B| / \sqrt{|A||B|})$. Der Kernschritt
zeigt: Ist $d(A, A)$ klein, so ist $A$ nahe (im Ruzsa-Abstand) an einem Unterraum.
Die Entropiemethode, kombiniert mit dem \emph{Faserungstrick} und einem
Dimensions-Reduktions-Argument, liefert die polynomiale Schranke.
\end{proof}

\begin{vermutung}[Optimale PFR-Konstante]\label{verm:pfr_optimal}
Der optimale Exponent in Satz~\ref{satz:pfr_f2} ist $\leq 6$ (statt der bewiesenen
$14$).
\end{vermutung}

%% ============================================================
\section{Quantitative Schranken: Übersichtstabelle}\label{sec:tabelle}
%% ============================================================

\begin{center}
\begin{tabular}{lll}
\toprule
\textbf{Ergebnis} & \textbf{Schranke für $|P|/|A|$} & \textbf{Dimension $d$} \\
\midrule
Freiman (1966)   & $\exp(\exp(K))$        & $\exp(K)$ \\
Ruzsa (1994)     & $\exp(K^4 \log K)$     & $2K^4 \log(2K)$ \\
Bilu (1999)      & $\exp(O(K^2))$         & $O(K^2)$ \\
Chang (2002)     & $\exp(O(K^2))$         & $O(K)$ \\
Sanders (2012)   & $\exp(O(\log^4 K))$    & $O(\log K)$ \\
PFR über $\F_2^n$ (2023) & $K^{14}$ (polynomial) & --- \\
\bottomrule
\end{tabular}
\end{center}

%% ============================================================
\section{Offene Probleme}\label{sec:offen}
%% ============================================================

Wir führen die wichtigsten offenen Probleme des Gebiets auf.

\begin{vermutung}[Polynomiales Freiman--Ruzsa über $\Z$]\label{verm:pfr_z}
Es gibt absolute Konstanten $c, C > 0$ derart, dass für $A \subset \Z$ mit
$|A+A| \leq K|A|$ die Menge $A$ in einer VAP $P$ mit $|P| \leq K^C |A|$ und
Dimension $d \leq c \log K$ enthalten ist.
\end{vermutung}

\begin{vermutung}[Optimale Verdoppelungskonstante]\label{verm:opt_verd}
Arithmetische Progressionen $P$ der Länge $n$ erfüllen $\sigma[P] = 2 - 1/|P|$.
Es wird vermutet, dass Mengen mit $\sigma[A] = 2 - 1/|A| + \varepsilon$ für kleines
$\varepsilon$ stets "nahe" an arithmetischen Progressionen liegen, aber ein
präziser Stabilitätssatz mit optimalen Konstanten ist offen.
\end{vermutung}

\begin{vermutung}[Freiman--Bilu-Exponent]\label{verm:fb_exp}
Der korrekte Exponent $\alpha$ in $|P| \leq K^\alpha |A|$ in
Vermutung~\ref{verm:pfr_z} erfüllt $1 \leq \alpha \leq 2$.
\end{vermutung}

\begin{vermutung}[Inverses Summenmengenproblem für nichtabelsche Gruppen]
\label{verm:nichtabelsch}
Ein geeignetes Analogon von Freimanns Satz gilt für nichtabelsche Gruppen
(mit der Rolle von VAPs besetzt durch Nebenklassen approximativer Gruppen im
Sinne von Hrushovski und Breuillard--Green--Tao).
\end{vermutung}

\begin{bemerkung}
Breuillard, Green und Tao \cite{BreuillardGreenTao2012} bewiesen eine Version von
Freimanns Satz für lineare Gruppen über Körpern: Mengen mit kleiner Verdoppelung
in $GL_n(k)$ sind im Wesentlichen "nilpotente Nebenklassen-Progressionen".
\end{bemerkung}

%% ============================================================
\section{Schluss}
%% ============================================================

Freimanns Satz steht an der Schnittstelle von additiver Kombinatorik, harmonischer
Analysis und Ergodentheorie. Ausgehend von der qualitativen Einsicht, dass kleine
Verdoppelung arithmetische Struktur erzwingt, hat die quantitative Theorie enorme
Fortschritte gemacht: von Ruzsa's exponentieller Schranke (1994) über Sanders'
quasi-polynomiale Verbesserung (2012) bis zu Gowers--Green--Manners--Tao's
polynomialem Ergebnis über $\F_2^n$ (2023). Das zentrale offene Problem —
ob eine polynomiale Schranke über $\Z$ gilt — bleibt eine der herausforderndsten
und wichtigsten Fragen des Gebiets.

%% Literaturverzeichnis
\begin{thebibliography}{99}

\bibitem{Bogolyubov1939}
N.\,N.~Bogolyubov,
\textit{On some arithmetical properties of almost periodic functions},
Ann.\ Chaire Math.\ Phys.\ Kiev \textbf{4} (1939), 185--205.

\bibitem{Bilu1999}
Y.~Bilu,
\textit{Structure of sets with small sumset},
Ast\'erisque \textbf{258} (1999), 77--108.

\bibitem{BreuillardGreenTao2012}
E.~Breuillard, B.~Green, T.~Tao,
\textit{The structure of approximate groups},
Publ.\ Math.\ IHES \textbf{116} (2012), 115--221.

\bibitem{Freiman1966}
G.\,A.~Freiman,
\textit{Foundations of a Structural Theory of Set Addition},
Kazan' Gos.\ Ped.\ Inst., Kazan', 1966 (russisch);
engl.\ Übers.\ AMS, Providence, 1973.

\bibitem{GreenRuzsa2007}
B.~Green, I.\,Z.~Ruzsa,
\textit{Freiman's theorem in an arbitrary abelian group},
J.\ Lond.\ Math.\ Soc.\ \textbf{75} (2007), 163--175.

\bibitem{GowersTao2023}
W.\,T.~Gowers, B.~Green, F.~Manners, T.~Tao,
\textit{On a conjecture of Marton},
Preprint arXiv:2311.05762 (2023).

\bibitem{Plunnecke1970}
H.~Pl\"unnecke,
\textit{Eine zahlentheoretische Anwendung der Graphentheorie},
J.\ reine angew.\ Math.\ \textbf{243} (1970), 171--183.

\bibitem{Roth1953}
K.\,F.~Roth,
\textit{On certain sets of integers},
J.\ Lond.\ Math.\ Soc.\ \textbf{28} (1953), 104--109.

\bibitem{Ruzsa1989}
I.\,Z.~Ruzsa,
\textit{An application of graph theory to additive number theory},
Scientia (Budapest) \textbf{3} (1989), 97--109.

\bibitem{Ruzsa1994}
I.\,Z.~Ruzsa,
\textit{Generalized arithmetical progressions and sumsets},
Acta Math.\ Hungar.\ \textbf{65} (1994), 379--388.

\bibitem{Ruzsa2009}
I.\,Z.~Ruzsa,
\textit{Sumsets and structure},
in: Combinatorial Number Theory and Additive Group Theory,
Birkh\"auser, 2009, 87--210.

\bibitem{Sanders2012}
T.~Sanders,
\textit{On the Bogolyubov--Ruzsa lemma},
Analysis \& PDE \textbf{5} (2012), 627--655.

\bibitem{Schur1916}
I.~Schur,
\textit{\"Uber die Kongruenz $x^m + y^m \equiv z^m \pmod{p}$},
Jahresber.\ Dtsch.\ Math.-Ver.\ \textbf{25} (1916), 114--117.

\bibitem{TaoVu2006}
T.~Tao, V.~Vu,
\textit{Additive Combinatorics},
Cambridge University Press, Cambridge, 2006.

\end{thebibliography}

\end{document}
